\section{满足RLEB的例子}
\label{sec:rank-deficient-hamilton}

考虑秩亏设计矩阵 $(1,0)$ 对应的幂次回归损失
\[
 f(x)=\frac23|x_1|^{3/2},\qquad x=(x_1,x_2)\in\mathbb R^2,
\]
以及动能 $k(p)=\frac23\|p\|^{3/2}$。
Maddison 等在 \emph{Hamiltonian Descent Methods} 中研究的
全隐式 Hamilton 更新为
\begin{equation}
\label{eq:hd-implicit}
 x^{j+1}-x^j=\nabla k(p^{j+1}),\qquad
 p^{j+1}-p^j
 =-\nabla f(x^{j+1})-\Gamma p^{j+1}.
\end{equation}
下述回归实例取 $\Gamma=99$。记
\[
 g(x)=\nabla f(x)
   =\bigl(\operatorname{sgn}(x_1)\sqrt{|x_1|},0\bigr),
 \qquad
 h(p)=\nabla k(p)
   =
   \begin{cases}
     p/\sqrt{\|p\|},&p\ne0,\\
     0,&p=0.
   \end{cases}
\]
在 $X=\mathbb R^2\times\mathbb R^2$ 上定义
\[
 F(x,p)=\bigl(-h(p),\,g(x)+99p\bigr).
\]
于是 \eqref{eq:hd-implicit} 等价于
$z^{j+1}=J_F(z^j)$，其中 $z^j=(x^j,p^j)$。零点集为
\[
 S=F^{-1}(0)
   =\{((0,s),0):s\in\mathbb R\}.
\]

\begin{proposition}
\label{prop:rank-deficient-hamilton-rleb}
令 $U=X$、$R=10^{-6}$，并将 $F$ 的图像限制在
近端输入属于
$U_R=\{z:d(z,S)\leq R\}$ 的部分。
该图像块满足 (A1)--(A4)，可取
\[
 \lambda=1,\quad \gamma=\frac12,\quad L=\frac12,\quad
 \bar t=1,\quad \psi(t)=3t^2,\quad \kappa=\frac14.
\]
\end{proposition}

\begin{proof}
将近端输入记为 $(a,b)$、输出记为 $(x,p)$。
近端方程等价于
\begin{equation}
\label{eq:hd-proximal-equations}
 a=x-h(p),\qquad b=g(x)+100p.
\end{equation}
由于 $k^*(u)=\|u\|^3/3$，$x$ 是
\[
 u\longmapsto
 k^*(u-a)+\frac1{100}f(u)-\frac1{100}\langle b,u\rangle
\]
的唯一极小点，且 $p=\nabla k^*(x-a)$。
故 $J_F$ 在 $X$ 上单值且处处有定义。
对任意 $s\in S$，$(s,0)$ 属于所取图像块；
这验证了 (A1)。

以下两个估计用于检验其余条件。
首先，对 $(a,b)\in U_R$，由
$|a_1|\leq R$、$\|b\|\leq R$ 及
\eqref{eq:hd-proximal-equations} 得
\begin{equation}
\label{eq:hd-momentum-small}
 \|p\|\leq R.
\end{equation}
事实上，若 $p_1g_1(x)\geq0$，则
$100\|p\|\leq\|b\|$；否则 $x_1$ 与 $p_1$ 异号，
因而
$|p_1|\leq R\sqrt{\|p\|}$、
$|p_2|\leq R/100$，仍推出
\eqref{eq:hd-momentum-small}。

其次，\eqref{eq:hd-proximal-equations} 给出
$p_2=b_2/100$，而 $p_1$ 由严格递增的方程
\[
 b_1=100p_1+
 g_1\bigl(a_1+h_1(p_1,b_2/100)\bigr)
\]
确定。在可微点作隐式微分，并使用
\[
 \frac{\partial h_1}{\partial p_1}
       \geq\frac1{2\sqrt{\|p\|}},
 \qquad
 \left|
 \frac{\partial h_1/\partial p_2}
      {\partial h_1/\partial p_1}
 \right|\leq\frac{\sqrt2}{4},
\]
可得输入到输出动量的 Jacobian 范数小于 $0.02$。
对 $g_1,h$ 作平滑正则化，上述界一致成立；
沿凸集 $U_R$ 中的线段积分并取极限，得到
\begin{equation}
\label{eq:hd-momentum-lipschitz}
 \|p(a,b)-p(\tilde a,\tilde b)\|
 \leq0.02\|(a,b)-(\tilde a,\tilde b)\|.
\end{equation}

记 $C=2J_F-I$。利用
$\|h(p)-h(\tilde p)\|^2\leq2\|p-\tilde p\|$
及 \eqref{eq:hd-momentum-lipschitz}，当两个近端输入
$(a,b),(\tilde a,\tilde b)\in U_R$ 的距离为
$\delta\leq R$ 时，有
\[
 \begin{split}
 \|C(a,b)-C(\tilde a,\tilde b)\|
 &\leq 2\sqrt{0.04\,\delta}+1.04\,\delta\\
 &\leq0.40104\sqrt{\delta}
 <L\delta^{1/2}.
 \end{split}
\]
因此 (A2) 对所取图像块上的任意点对成立。

设 $(x,p)=J_F(a,b)$、$(a,b)\in U_R$，并记
$\rho=r_F(x,p)=\|F(x,p)\|$。由
\[
 \rho^2=\|p\|+\|g(x)+99p\|^2
\]
和 \eqref{eq:hd-momentum-small} 得
\[
 \|p\|\leq\rho^2,\qquad
 \|g(x)\|\leq(1+99\sqrt R)\rho.
\]
由于 $|x_1|=\|g(x)\|^2$，故
\[
 d((x,p),S)
 \leq |x_1|+\|p\|
 \leq\bigl((1+99\sqrt R)^2+1\bigr)\rho^2
 <3\rho^2.
\]
又有 $(R+L\sqrt R)/2<1=\bar t$，
从而 (A3) 成立。最后，对 $0<r\leq R$，
\[
 \frac1r\,
 \psi\!\left(\frac{r+L\sqrt r}{2}\right)
 =\frac34\left(\sqrt r+\frac12\right)^2
 \leq\frac34(0.001+0.5)^2
 <\frac14.
\]
这验证了 (A4)。
\end{proof}

由主定理，对每个满足
$d_0:=d(z^0,S)\leq10^{-6}$ 的初值，
隐式迭代 \eqref{eq:hd-implicit} 收敛到某个
$z^\infty\in S$，并且
\[
 d(z^j,S)\leq4^{-j}d_0,\qquad
 \|z^j-z^\infty\|
 \leq\frac23\,4^{-j}d_0
      +\frac12\,2^{-j}\sqrt{d_0}.
\]
该结论允许非零初始动量，并确定整条位置--动量
轨道的极限点，而不只控制目标函数值。

本例也给出与已有条件的具体区别。
Maddison 等的隐式法收敛命题所用能量配对条件
要求以统一常数控制
$|\langle\nabla f(x),\nabla k(p)\rangle|$
与 $f(x)+k(p)$ 的比值。但取
$x=p=(t,0)$、$t\downarrow0$，该比值等于
\[
 \frac{t}{\frac43t^{3/2}}
 =\frac{3}{4\sqrt t}\longrightarrow\infty.
\]
因此不能由该命题推出上述收敛结论。

Luke--Tam 命题~4 对应的局部次单调性条件则要求
反射预解式在所考虑的完整局部图像上满足
Lipschitz 估计。令
$w_t=((0,0),(0,t))$、$w_0=0$。由
\eqref{eq:hd-proximal-equations}，
\[
 J_F(w_t)
 =\bigl((0,\sqrt{t/100}),(0,t/100)\bigr),
\]
故
\[
 \frac{\|(2J_F-I)w_t-(2J_F-I)w_0\|}
      {\|w_t-w_0\|}
 \geq\frac1{5\sqrt t}\longrightarrow\infty.
\]
所比较的点均趋于零点并属于上述图像块，
故该局部 Lipschitz 条件在这里不成立。